\documentclass[sigconf]{aamas}

\usepackage{balance}
\usepackage{booktabs}
\usepackage{graphicx}
\usepackage{amsmath}
\usepackage{multirow}

\setcopyright{ifaamas}
\acmConference[AAMAS '23]{Proc.\@ of the 22nd International Conference
on Autonomous Agents and Multiagent Systems (AAMAS 2023)}{May 29 -- June 2, 2023}
{London, United Kingdom}{A.~Ricci, W.~Yeoh, N.~Agmon, B.~An (eds.)}
\copyrightyear{2023}
\acmYear{2023}
\acmDOI{}
\acmPrice{}
\acmISBN{}

\acmSubmissionID{???}

\title{Multiwinner Voting for Evidence Selection in Knowledge-Based Visual Question Answering}

\author{Your Name}
\affiliation{
  \institution{Carleton University}
  \city{Ottawa}
  \country{Canada}}
\email{your.email@carleton.ca}

\begin{abstract}
% ~150 words. Hit these beats:
% 1. Problem: KB-VQA needs external evidence, top-k retrieval picks redundant passages
% 2. Idea: recast evidence selection as multiwinner election — LLM agents vote, voting rules select
% 3. What you tested: AV, PAV, CC, MES, Borda on OK-VQA with frozen LLaVA
% 4. Key result: Borda +2.8%, CC +2.6% over top-k
% 5. Insight: coverage (CC) matters more than proportionality (PAV/MES) with few agents
\end{abstract}

\keywords{Multiwinner Voting, Visual Question Answering, Multi-Agent Systems, Evidence Selection}

\begin{document}

\pagestyle{fancy}
\fancyhead{}

\maketitle

%% ============================================================
\section{Introduction}
%% ============================================================

% ~1 page. Structure:
%
% Para 1: KB-VQA requires external knowledge. Current approach is top-k retrieval.
%          Problem: top-k picks redundant passages (all say the same thing).
%
% Para 2: Key insight — evidence selection IS a committee election.
%          Passages = candidates, agents = voters, selected evidence = committee.
%          This isn't just an analogy — voting theory gives us formal guarantees
%          (JR, EJR) that top-k and MMR lack.
%
% Para 3: What you did — 3 LLM agents with different perspectives vote on passages,
%          5 voting rules aggregate their ballots. Tested on OK-VQA with frozen LLaVA.
%
% Para 4: Contributions as a numbered list:
%          1. Borda/CC beat top-k by ~2.7%
%          2. Coverage (CC) matters more than proportionality (PAV/MES collapse to AV)
%          3. Advantage is largest at k=5, shrinks at large k

%% ============================================================
\section{Related Work}
%% ============================================================

% ~0.5-1 page, 3 paragraphs. This goes EARLY in AAMAS papers — position your work
% relative to existing literature before diving into definitions.
%
% Para 1: Retrieval-augmented VQA — RA-VQA, KRISP, Prophet.
%          They focus on WHAT to retrieve (better retrievers, better prompts).
%          You focus on HOW TO SELECT from retrieved candidates — complementary contribution.
%
% Para 2: Diversity in retrieval — MMR, DPP (determinantal point processes).
%          These are heuristic/probabilistic. Your approach is principled (voting theory)
%          with formal guarantees (JR, EJR). MMR is actually one of your baselines.
%
% Para 3: Multiwinner voting in AI — Skowron et al. on data summarization,
%          Lackner & Skowron survey on approval-based committee voting.
%          You apply these ideas to a new domain: multimodal evidence selection.
%          Nobody has connected multiwinner voting to VQA before — that's your novelty.

%% ============================================================
\section{Preliminaries}
%% ============================================================

% Formal definitions the reader needs before understanding your method.

\subsection{Knowledge-Based VQA}
% 1 paragraph: What is KB-VQA, why does it need retrieval, cite OK-VQA.
% Keep it short — reader just needs to know the task setup.

\subsection{Multiwinner Voting}
% Define the problem: n voters, m candidates, select committee of size k.
% Then define each rule in 1-2 sentences each:
%   - AV: most approvals win
%   - PAV: diminishing returns (1/j marginal gain), greedy
%   - CC: greedy max coverage of voters
%   - MES: equal budgets, candidates purchased by supporters
%   - Borda: rank-based scoring (m - rank)
% Cite Lackner & Skowron 2023 for the survey.

\subsection{Proportional Representation}
% Define JR and EJR in plain language (1 paragraph):
%   - JR: if n/k voters share an approved candidate, they deserve 1 representative
%   - EJR: stronger — l*n/k cohesive voters deserve l representatives
%   - Which rules satisfy what: PAV -> JR, MES -> EJR, CC -> coverage, AV/Borda -> neither
% This sets up the later finding that JR/EJR don't matter with 3 agents.

%% ============================================================
\section{Method}
%% ============================================================

\subsection{Pipeline Overview}
% Draw the pipeline figure (or use ASCII art):
%   Image+Question -> BLIP caption -> BM25 top-20 -> 3 agents -> voting rule -> LLaVA -> answer
% Describe each stage in 1 sentence.

\subsection{Agent Design}
% Describe the 3 agent roles and WHY you chose these perspectives:
%   - Relevance: broad topical filter
%   - Uniqueness: promotes diversity (non-redundant passages)
%   - Gap-filler: catches what obvious relevance misses
% Mention: each outputs APPROVED set + RANKING (enables both approval and Borda)
% Mention: prompted to approve >= 5 passages (avg ~6.2 approvals)

\subsection{Voting Rules}
% Give formal definitions for each rule applied to your setting.
% Mention tie-breaking with BM25 scores.
% Define the two baselines: top-k (BM25 scores only) and MMR (relevance-diversity tradeoff).

%% ============================================================
\section{Experimental Setup}
%% ============================================================

% ~0.5 page. Cover:
% - Dataset: OK-VQA, 500 questions, standard VQA accuracy metric (min(1, matches/3))
% - Knowledge base: 3M Wikipedia passages, BM25 index, top-20 candidates
% - Models (all frozen, no fine-tuning):
%     - BLIP-base (captioner)
%     - Mistral-7B-Instruct-v0.3, 4-bit quantized (agents)
%     - LLaVA-1.5-13B, 4-bit quantized (answerer)
% - Committee size k=5, 3 agents
% - Hardware: 4x NVIDIA L4 GPUs

%% ============================================================
\section{Results}
%% ============================================================

\subsection{Main Results}
% Table: all 7 methods, accuracy, delta vs top-k
% Key points to make:
%   - All voting rules beat top-k
%   - Borda (0.486) and CC (0.484) are the clear winners
%   - Even simple AV (+1.1%) helps — multi-agent assessment adds value
%   - MMR (+1.5%) is decent but voting rules do better

\subsection{Committee Size Ablation}
% Table/figure: k = 3, 5, 7, 10 for each rule
% Points:
%   - Smaller k is universally better (less noise for VLM)
%   - CC's advantage is largest at k=5 (+2.6%), narrows at k=3 and k=10
%   - At k=3: all methods converge (too few slots for selection to differ)
%   - At k=10: half the pool selected regardless, differences wash out

\subsection{Agent Count Ablation}
% Table/figure: N = 1, 2, 3 agents
% Points:
%   - N=1: voting is meaningless, just one agent's opinion
%   - N=2: some diversity, should see improvement
%   - N=3: full system, best results
%   - This proves multi-agent DISAGREEMENT is what drives the improvement
%%% TODO: fill in after agent ablation experiment completes

\subsection{Proportionality Analysis}
% Key findings:
%   - AV = PAV = MES (identical committees, Jaccard 1.0) — proportionality collapses with 3 agents
%   - Zero JR/EJR violations across all rules — ballots are large enough that every agent
%     always gets represented trivially
%   - This means: proportionality is a freebie, but coverage (CC) is the actual differentiator
%   - Implication: with few agents, use CC; proportionality would matter with more agents

\subsection{Committee Overlap}
% Jaccard similarity table between all pairs of rules
% Key point: CC has lowest overlap with top-k (0.12) — selects most different passages
% CC's unique passage selections are the ones that improve accuracy

\subsection{Qualitative Examples}
% Show 2-3 cases where CC gets it right and top-k doesn't
% For each: question topic, top-k answer (wrong), CC answer (right), ground truth
% Explain WHY: CC picked a passage covering a different angle

%% ============================================================
\section{Discussion}
%% ============================================================

% ~0.5 page. Cover:
% - Why coverage > proportionality at this scale (3 agents too few for coalitions)
% - Why smaller k helps (VLMs are noise-sensitive in context)
% - Limitations:
%     - BM25 retrieval (not dense retrieval)
%     - Single dataset (OK-VQA)
%     - Single VLM (LLaVA-1.5-13B)
%     - 3 agents (doesn't exercise proportionality)
% - Future work: more agents, dense retrieval, multiple datasets, adaptive k

\balance

%% ============================================================
\section{Conclusion}
%% ============================================================

% ~0.5 page. Three beats:
% 1. What you did: recast evidence selection as multiwinner election
% 2. What you found: CC and Borda beat top-k; coverage > proportionality at this scale
% 3. Why it matters: principled, formal alternative to heuristic diversity methods

\begin{acks}
%%% TODO: Acknowledge supervisor, course (COMP 5900), compute resources
\end{acks}

\bibliographystyle{ACM-Reference-Format}
\bibliography{references}

\end{document}
